\section{The error of the truncated expansion at every temperature}
\label{sec:hastings_temperature_error}

At an arbitrary inverse temperature $\beta$, the truncated operator is used in
$n$ factors: Hastings replaces $e^{-\beta H}=(e^{-\beta H/n})^n$ by the product
$\rho(\beta/n,l)^n$ of $n$ truncated factors, each at the inverse temperature
$\beta/n$ \cite[lines 1303--1314]{Hastings2006Solving}. The error is the sum of
the operators $\rho_m$ of the terms with at least $m$ large clusters across the
$n$ factors \cite[lines 1334--1341]{Hastings2006Solving}, and it is estimated
``following the same arguments as before''
\cite[lines 1342--1347]{Hastings2006Solving}. This section carries out that
estimate. Throughout, $y=(e^{J\beta/n}-1)e^{SJ\beta/n}$ and
$Z=\operatorname{tr}e^{-\beta H}$.

\begin{theorem}[The error of the product of truncated factors]
    \label{thm:hastings_temperature_error}
    \lean{TNLean.Lattice.ClusterExpansion.sum_prod_le_exp_of_subset_powerset,
        TNLean.Lattice.ClusterExpansion.norm_trace_mul_exp_sub_truncatedExp_pow_le}
    \leanok
    \uses{thm:hastings_truncation_inclusion_exclusion,
        thm:hastings_cluster_part_product, thm:hastings_percolation_exp}
    Let $H=\sum_kH_k$ be Hermitian with Hermitian terms of norm at most $J$, let
    $\beta\ge0$ and $n\ge1$, and let every large connected set $C$ have
    $n_C+n_{\partial C}\le Sn_C$. Then for every operator $O$
    \begin{align}
        \bigl|\operatorname{tr}\bigl(O(e^{-\beta H}-\rho(\beta/n,l)^n)\bigr)\bigr|
        \le\|O\|\Bigl(\exp\Bigl(n\sum_Cy^{n_C}\Bigr)-1\Bigr)Z,
        \label{eq:hastings_temperature_error}
    \end{align}
    summed over the large connected sets $C$.
\end{theorem}

\begin{proof}
    \leanok
    \uses{thm:hastings_truncation_inclusion_exclusion,
        thm:hastings_words_inclusion_exclusion, thm:hastings_cluster_part_product,
        thm:hastings_percolation_exp, thm:hastings_clusters_meet}
    By Theorem~\ref{thm:hastings_truncation_inclusion_exclusion} at the inverse
    temperature $\beta/n$, each factor is $\sum_U(-1)^{|U|}\rho(\beta/n,U)$ over
    the admissible families $U$ of large connected sets, and the empty family
    gives $e^{-\beta H/n}$. Expanding the $n$-th power gives a sum over the
    sequences $(U_0,\dots,U_{n-1})$ of admissible families of the ordered
    products $\prod_j(-1)^{|U_j|}\rho(\beta/n,U_j)$, and the sequence of empty
    families gives $e^{-\beta H}$. The members of an admissible family are
    pairwise disjoint (Theorem~\ref{thm:hastings_clusters_meet}), so the union
    of $U_j$ has $\sum_{C\in U_j}n_C$ terms and satisfies the hypothesis of
    Theorem~\ref{thm:hastings_cluster_part_product}, which bounds the term of
    every other sequence by $\|O\|Z\prod_jf(U_j)$ with
    $f(U)=\prod_{C\in U}y^{n_C}$. Summing over the sequences other than the
    empty one gives $\|O\|Z\bigl((\sum_Uf(U))^n-1\bigr)$. The nonempty admissible
    families are among the nonempty sets of large connected sets, so
    Theorem~\ref{thm:hastings_percolation_exp} gives
    $\sum_Uf(U)\le\exp(\sum_Cy^{n_C})$, and the $n$-th power of this is the
    exponential in~(\ref{eq:hastings_temperature_error}).
\end{proof}

Grouping the sequences by their total number $m$ of large clusters recovers the
bound $(1/m!)(n\sum_Cy^{n_C})^m$ on the operator $\rho_m$
of~\cite[lines 1334--1341]{Hastings2006Solving}.

\begin{theorem}[The error at every temperature, percolation form]
    \label{thm:hastings_temperature_error_percolation}
    \lean{TNLean.Lattice.ClusterExpansion.norm_trace_mul_exp_sub_truncatedExp_pow_le_percolation}
    \leanok
    \uses{thm:hastings_temperature_error, thm:hastings_percolation_cluster_sum,
        thm:hastings_percolation}
    In the setting of Theorem~\ref{thm:hastings_temperature_error}, let every
    large connected set have at least $l+1$ terms, let at most $Nx^k$ large
    connected sets have $k$ terms, for every $k$, where $N,x\ge0$, and let
    $0<xy\le1/2$. Put $\xi=1/\ln(1/(xy))$. Then for every operator $O$
    \begin{align*}
        \bigl|\operatorname{tr}\bigl(O(e^{-\beta H}-\rho(\beta/n,l)^n)\bigr)\bigr|
        \le\|O\|\bigl(\exp(nNe^{-l/\xi})-1\bigr)Z.
    \end{align*}
    The source prints the bound
    $\{\exp[(\beta/\beta_0)Ne^{-l/\xi}]-1\}\|O\|$ on
    $Z^{-1}|\operatorname{tr}O\rho(\beta)-\operatorname{tr}O\rho(\beta,l)|$
    \cite[lines 1342--1347]{Hastings2006Solving}, with $\beta/\beta_0$ in place of
    the number $n$ of factors. Its $n$, the smallest even integer with
    $\beta/n\le\beta_0$, satisfies $\beta/\beta_0\le n<\beta/\beta_0+2$, and its bound
    on $\rho_m$ gives the exponent with $n$, which is the bound proved here
    \cite{gap:hastings05_temperature_factor_count}. The threshold $l$ counts the
    terms of a large connected set rather than its diameter, as in
    Theorem~\ref{thm:hastings_percolation}
    \cite{gap:hastings05_percolation_cluster_size}.
\end{theorem}

\begin{proof}
    \leanok
    \uses{thm:hastings_temperature_error, thm:hastings_percolation_cluster_sum,
        thm:hastings_percolation}
    Theorem~\ref{thm:hastings_percolation_cluster_sum} bounds
    $\sum_Cy^{n_C}$ by $N(xy)^{l+1}/(1-xy)$, which is at most
    $N(xy)^l=Ne^{-l/\xi}$ since $xy\le1/2$, as in
    Theorem~\ref{thm:hastings_percolation}. Insert this into
    Theorem~\ref{thm:hastings_temperature_error}.
\end{proof}
